% © 2026 Ing. David Jaroš — CC BY-NC-ND 4.0
%
% This work is licensed under a Creative Commons Attribution-NonCommercial-NoDerivatives
% 4.0 International License (CC BY-NC-ND 4.0).
% See LICENSE.md for full license text.

% UBT-AI-PROVENANCE-BEGIN
% schema: ubt-ai-provenance/v1
% tier: C_working
% ai_assistance: disclosed
% human_review: risk-based
% editorial_responsibility: Ing. David Jaroš
% policy: ../../AI_PROVENANCE.md
% notice: Working material; exhaustive human review is not claimed.
% UBT-AI-PROVENANCE-END

\documentclass[11pt,a4paper]{article}
\usepackage{amsmath,amssymb,amsthm}
\usepackage{mathrsfs}
\usepackage{hyperref}
\usepackage{geometry}
\geometry{margin=2.5cm}

\newtheorem{proposition}{Proposition}
\newtheorem{lemma}{Lemma}
\newtheorem{remark}{Remark}
\newtheorem{gap}{Open Gap}

\newcommand{\proved}[1]{\textbf{[PROVED]} #1}
\newcommand{\heuristic}[1]{\textbf{[HEURISTIC]} #1}
\newcommand{\speculative}[1]{\textbf{[SPECULATIVE]} #1}

\title{Loop Integrals Appendix: One-Loop Vacuum Polarisation in $d = 1$}
\author{Ing.\ David Jaro\v{s}\\
\small \texttt{research\_tracks/rg\_nlogn/}}
\date{May 2026}

\begin{document}
\maketitle

\begin{abstract}
This appendix collects explicit one-loop loop integral calculations in
$d = 1$ relevant to the UBT KK effective potential $V_\mathrm{eff}(n) = n^2 - Bn\ln n$.
We perform the dimensional regularisation calculation in $d = 1 - \epsilon$,
identify UV-divergent and UV-finite parts, and derive the $n^2\ln n$ coefficient.
We also show the two-loop sunset diagram estimate and explain the normalisation convention.
\end{abstract}

%------------------------------------------------------------
\section{One-Loop Integral in $d = 1$}
%------------------------------------------------------------

\subsection{Bubble diagram}

The relevant one-loop diagram is the bubble (vacuum polarisation) with one
massive propagator of mass $m$:

\[
  I_1(m^2) = \int \frac{d^dk}{(2\pi)^d} \frac{1}{k^2 + m^2}
\]

where $d = 1 - \epsilon$ (Euclidean signature, Wick-rotated).

Using the standard formula for Euclidean loop integrals:
\[
  \int \frac{d^dk}{(2\pi)^d} \frac{1}{(k^2 + m^2)^n}
  = \frac{1}{(4\pi)^{d/2}} \frac{\Gamma(n - d/2)}{\Gamma(n)}
    \frac{1}{(m^2)^{n-d/2}},
\]

for $n = 1$, $d = 1 - \epsilon$:
\[
  I_1(m^2) = \frac{1}{(4\pi)^{(1-\epsilon)/2}}
    \Gamma\!\left(\tfrac{1+\epsilon}{2}\right) (m^2)^{-(1+\epsilon)/2}.
\]

Expanding in $\epsilon$:
\[
  I_1(m^2) = \frac{1}{\sqrt{4\pi}}
  \Gamma\!\left(\tfrac12\right)(m^2)^{-1/2}
  \left[1 - \frac{\epsilon}{2}\ln(4\pi m^2) + \ldots\right].
\]

Since $\Gamma(\frac12) = \sqrt{\pi}$:
\[
  I_1(m^2) = \frac{1}{2\sqrt{\pi}} \cdot \sqrt{\pi} \cdot m^{-1}
  + \mathcal{O}(\epsilon)
  = \frac{1}{2m} + \mathcal{O}(\epsilon).
\]

\textbf{Proof status}: \proved{Standard dimensional regularisation; see
\cite[App.\ A]{Peskin}.}

\subsection{One-loop effective potential from bubble diagram}

The one-loop contribution to the effective potential from a field of mass $m$
in $d = 1$ (i.e.\ a quantum mechanical system on a line/circle) is

\[
  V_\mathrm{1-loop}(m^2)
  = -\frac{1}{2}\int \frac{dk}{2\pi} \ln(k^2 + m^2)
  = -\frac{1}{2} \cdot \frac{m}{2\pi} \cdot 2\pi \cdot \frac{\partial}{\partial m^2}
    \int_0^\infty \frac{dk}{\pi} \frac{1}{k^2 + m^2}
\]

More directly, using the heat-kernel representation:
\[
  \ln(k^2 + m^2) = -\int_0^\infty \frac{dt}{t} e^{-t(k^2 + m^2)},
\]

so (in $d = 1$ with momentum measure $dk/(2\pi)$):
\begin{align*}
  V_\mathrm{1-loop}(m^2)
  &= \frac{1}{2}\int_0^\infty \frac{dt}{t}\, e^{-tm^2}
     \int_{-\infty}^\infty \frac{dk}{2\pi}\, e^{-tk^2} \\
  &= \frac{1}{2}\int_0^\infty \frac{dt}{t}\, e^{-tm^2}
     \frac{1}{\sqrt{4\pi t}} \\
  &= \frac{1}{4\sqrt{\pi}}\int_0^\infty t^{-3/2} e^{-tm^2}\, dt \\
  &= \frac{1}{4\sqrt{\pi}} \cdot \frac{2\sqrt{\pi}}{m}
   = \frac{1}{2m}.
\end{align*}

Wait — this gives $V_\mathrm{1-loop} = m/2$, the zero-point energy.  For a
massive field on a circle (KK mode), the correct one-loop vacuum energy density is:

\begin{proposition}[One-loop energy per KK mode, $d = 1$]
  For a massive scalar field of mass $m$ on $S^1_R$ with periodic boundary conditions,
  the regulated one-loop vacuum energy density in the $\overline{\mathrm{MS}}$ scheme is
  \[
    \mathcal{E}_n = \frac{m_n}{2} - \frac{m_n}{2}
    - \frac{m_n^2}{4\pi}\left(\ln\frac{m_n^2}{\mu^2} - 1\right)
    + \mathcal{O}(\epsilon)
  \]
  after subtracting the zero-point divergence by normal ordering.
\end{proposition}

\textbf{Proof status}: \heuristic{The subtraction of the zero-point divergence
is a renormalisation choice; in $d = 1$ the UV divergence is only linear (not
$\Lambda^2$), so the prescription affects the coefficient.  The
$m^2\ln(m^2/\mu^2)$ term is scheme-independent.}

%------------------------------------------------------------
\section{Normalisation and Logarithmic Coefficient}
%------------------------------------------------------------

\subsection{Coleman--Weinberg potential in $d = 4$}

For comparison, in $d = 4$ the Coleman--Weinberg result is:
\[
  V_\mathrm{CW}(\phi) = \frac{m^4(\phi)}{64\pi^2}
  \left(\ln\frac{m^2(\phi)}{\mu^2} - \frac{3}{2}\right),
\]

where $m^2(\phi) = m^2 + \lambda\phi^2/2$ is field-dependent.

\subsection{$d = 1$ analogue}

In $d = 1$ the corresponding result for a scalar of mass $m$ is:
\[
  V_\mathrm{1-loop}^{(d=1)}(m^2) = -\frac{m^2}{4\pi}
  \left(\ln\frac{m^2}{\mu^2} - 1\right).
\]

Setting $m^2 = n^2$ (the $n$-th KK mode):
\[
  \delta V_n = -\frac{n^2}{4\pi}\left(2\ln n - \ln\mu^2 - 1\right).
\]

The $\ln\mu^2$ term renormalises the $n^2$ coefficient; the finite
$-n^2\ln n/(2\pi)$ remains.

\textbf{Proof status}: \proved{The $m^2\ln m^2$ coefficient in $d=1$ is
scheme-independent.}

%------------------------------------------------------------
\section{Winding Mode Contribution}
%------------------------------------------------------------

At the self-dual radius $R = 1$, winding modes have mass $m_\mathrm{wind,n}^2 = n^2$,
identical to KK modes.  Each level $n$ therefore has \emph{two} contributions
(KK + winding) at the same mass, doubling the one-loop correction:

\[
  \delta V_n^\mathrm{KK+wind} = 2 \times \left(-\frac{n^2}{4\pi}\right)
  \left(2\ln n - \ln\mu^2 - 1\right).
\]

After renormalisation, the finite part is
\[
  \delta V_n^\mathrm{finite} = -\frac{n^2\ln n}{\pi}.
\]

This gives $B_\mathrm{1-loop,KK+wind} = n/\pi \approx 43.6$ for $n = 137$.

\textbf{Proof status}: \heuristic{The winding/KK degeneracy at self-dual
radius is used; this is a string-theory result, not yet derived from UBT.}

%------------------------------------------------------------
\section{Two-Loop Sunset Diagram Estimate}
%------------------------------------------------------------

The two-loop sunset (setting sun) diagram contributes to the effective
potential at order $g^2$:

\[
  V_\mathrm{2-loop}(m^2) \sim \frac{g^2}{(4\pi)^2} m^4\ln(m^2/\mu^2)
  \quad (d=4)
\]

In $d = 1$:
\[
  V_\mathrm{2-loop}^{(d=1)}(m^2) \sim \frac{g^2}{16\pi^2} m^2\ln(m^2/\mu^2).
\]

For $m^2 = n^2$, this adds $\delta B_\mathrm{2-loop} \sim g^2 n/(16\pi^2)$.
For $n = 137$, $g \approx 0.3$:
\[
  \delta B_\mathrm{2-loop} \sim \frac{0.09}{250}\times 137 \approx 0.05.
\]

This is negligibly small — two-loop corrections cannot explain the
discrepancy $\Delta B \approx 2.4$ between 43.6 and 46.

\begin{gap}[Source of missing $\Delta B \approx 2.4$]
  The two-loop estimate contributes $\delta B \approx 0.05$, not $2.4$.
  Possible sources of the remaining discrepancy:
  \begin{enumerate}
    \item Gauge field loops: gluons, $W/Z$, $\gamma$ contribute with specific
          group-theory factors $b_0$.
    \item Fermion loops: contribute with opposite sign.
    \item Threshold corrections at the compactification scale.
    \item Non-perturbative corrections (instantons, etc.).
  \end{enumerate}
  \textbf{No derivation available.}
\end{gap}

%------------------------------------------------------------
\section{Dimensional Reduction Consistency}
%------------------------------------------------------------

\begin{proposition}[KK reduction consistency]
  The effective potential $V_\mathrm{eff}(n) = n^2 - Bn\ln n$ is consistent
  with a $(4+1)$-dimensional theory reduced on $S^1$, where the $(4+1)$-d
  mass term reduces to the KK tower $m_n^2 = n^2/R^2$.
\end{proposition}

\textbf{Proof status}: \proved{KK reduction on $S^1$ is standard; see
\cite{Kaluza1921,Klein1926}.}

\heuristic{The identification of the UBT compact direction with the fifth
Kaluza--Klein direction is motivated but not derived.}

%------------------------------------------------------------
\section{Summary}
%------------------------------------------------------------

\begin{center}
\begin{tabular}{ll}
\hline
Result & Status \\
\hline
$I_1(m^2) = 1/(2m)$ in $d=1$ & \textbf{Proved} \\
$V_\mathrm{1-loop} \propto m^2\ln m^2$ coefficient scheme-independent & \textbf{Proved} \\
$B_\mathrm{1-loop} \approx n/(2\pi)$ (one KK mode) & \textbf{Heuristic} \\
$B_\mathrm{KK+wind} \approx n/\pi$ (self-dual radius) & \textbf{Heuristic} \\
Two-loop: $\delta B \approx 0.05$ (negligible) & \textbf{Heuristic} \\
$B = 46$ exactly & \textbf{Open gap} \\
\hline
\end{tabular}
\end{center}

%------------------------------------------------------------
\begin{thebibliography}{9}
\bibitem{Peskin} M.E.\ Peskin and D.V.\ Schroeder, \emph{An Introduction to
  Quantum Field Theory}, Addison-Wesley, 1995.
\bibitem{Kaluza1921} T.\ Kaluza, ``Zum Unit\"atsproblem der Physik,''
  \emph{Sitzungsber.\ Preuss.\ Akad.\ Wiss.\ Berlin} (1921), 966--972.
\bibitem{Klein1926} O.\ Klein, ``Quantentheorie und f\"unfdimensionale
  Relativit\"atstheorie,'' \emph{Z.\ Phys.} \textbf{37} (1926), 895--906.
\bibitem{CW73} S.\ Coleman and E.\ Weinberg, ``Radiative corrections as the
  origin of spontaneous symmetry breaking,'' \emph{Phys.\ Rev.\ D} \textbf{7}
  (1973), 1888--1910.
\end{thebibliography}

\end{document}
